\documentclass{article}

\usepackage[utf8]{inputenc}
\usepackage[T1]{fontenc}
\usepackage[svgnames]{xcolor}
\usepackage{libertine}
\usepackage[frenchb]{babel}
% TODO \usepackage[french]{babelbib}
\usepackage{mathpartir}
\usepackage{syntax}
\usepackage{amssymb}
\usepackage{tikz}
\usepackage{amsmath}
\usepackage{float}
\usepackage{url}
\usepackage{xspace}
\usepackage{graphicx}
\usepackage[shortlabels]{enumitem}
\usepackage[backref=page]{hyperref}
\usepackage{links}
\usepackage[a4paper]{geometry}
\usepackage{amsthm}
\usepackage{stmaryrd}
\usepackage[normalem]{ulem}
\usepackage{wasysym}
\usepackage{multirow}
\usepackage{thmtools}
\usepackage{gramster}
\usepackage{switch}
\usepackage{pelagiccod}
\usepackage{conteq}

\newcommand{\irule}[3]{ \inferrule*[right=(#1)]{#2}{#3} }
\def\iaxiom#1#2{\irule{#1}{ }{#2}}

% http://tex.stackexchange.com/questions/3964/mathematical-symbol-for-disjoint-set-union
\makeatletter
\def\moverlay{\mathpalette\mov@rlay}
\def\mov@rlay#1#2{\leavevmode\vtop{%
   \baselineskip\z@skip \lineskiplimit-\maxdimen
   \ialign{\hfil$\m@th#1##$\hfil\cr#2\crcr}}}
\newcommand{\charfusion}[3][\mathord]{
    #1{\ifx#1\mathop\vphantom{#2}\fi
        \mathpalette\mov@rlay{#2\cr#3}
      }
    \ifx#1\mathop\expandafter\displaylimits\fi}
\makeatother

\newcommand{\dcup}{\charfusion[\mathbin]{\cup}{+}}
\newcommand{\bigdcup}{\charfusion[\mathbin]{\bigcup}{+}}

\def\opbin{\boxplus\xspace}
\def\opun{\boxminus\xspace}

% Types
\newcommand{\tInt}{\textsc{Int}\xspace}
\newcommand{\tFloat}{\textsc{Float}\xspace}
\newcommand{\tUnit}{\textsc{Unit}\xspace}

\def\structRec#1#2{\textsc{rec}(\&#1) #2} % TODO

\newcommand{\tFV}[1]{\textrm{FreeVars}(#1)}

\def\opxor{\mathbin{\char`\^}}

% Constant
\newcommand{\cNil}{\textsc{Nil}}

% Program & language names
\newcommand{\ctonewspeak}{\textsc{c2newspeak}}
\newcommand{\cpp}{\textsc{cpp}}
\newcommand{\newspeak}{\textsc{Newspeak}}

% Newspeak
\newcommand{\npkIf}[3]{\textsc{If}(#1) \{ #2 \} \textsc{Else} \{ #3 \}}
\newcommand{\npkDoWith}[2]{\textsc{Do}\{ #1 \}\textsc{With}~#2:}
\newcommand{\npkWhile}[1]{\textsc{While}(1)\{#1\}}
\newcommand{\npkGoto}[1]{\textsc{Goto}~#1}
\newcommand{\npkReturn}[1]{\textsc{Return}~#1}

% cML
\def\iPass{\textsc{Pass}\xspace}
\def\iIf#1#2#3{\npkIf{#1}{#2}{#3}}
\def\iWhile#1#2{\textsc{While}(#1)\{#2\}}
\def\iReturn#1{\textsc{Return}(#1)}
\def\iDecl#1#2#3{\textsc{Decl}~#1=#2~\textsc{in}\{#3\}}

\def\iThen#1#2{\textsc{If}(#1) \{ #2 \}}

\def\eNull{\textsc{Null}\xspace}
\def\eUnit{(\,)\xspace}

\def\vRet{\underline{R}}

% sets pour semantique
\def\sId{\textsc{Id}\xspace}
\def\sMem{\textsc{Mem}\xspace}
\def\sAddr{\textsc{Addr}\xspace}
\def\sInt{\textsc{Int}\xspace}
\def\sFloat{\textsc{Float}\xspace}
\def\sPhi{\Phi}

% Memory state (instructions)
% #1 : memory
% #2 : instruction
\def\msi#1#2{ \langle #2, #1 \rangle }

% Reduction (expressions)
\def\mm#1#2#3#4{ \msi{#1}{#2} → \msi{#3}{#4} }
\def\mms#1#2{\mm{m}{#1}{m}{#2}}

% Reduction (instructions)
\def\mmi#1#2#3#4{ \mm{#1}{#2}{#3}{#4} }

% Reduction (left-values)
\def\mmlv#1#2#3#4{ \mm{#1}{#2}{#3}{#4} }

% Typage sém
\def\semtyp#1#2#3{#1 ⊧ #2 : #3}

% Compatibilité Γ / m
\def\mcomp#1#2{#1 ⊧ #2}

% Reduction contexts
\def\ctxOp#1#2{#1~\opbin~#2}
\def\ctxEmpty{\bullet\xspace}
\def\ctxUnOp#1{\opun~#1}
\def\ctxLvDeref#1{*#1}
\def\ctxLvField#1#2{#1.#2}
\def\ctxLvIndex#1#2{#1[#2]}
\def\ctxSet#1#2{#1 ← #2}

% Context substitution
\def\ctxSub#1#2{#1\llparenthesis#2\rrparenthesis}
\def\ctxSubL#1#2{\ctxSub{#1}{#2}_L}
\def\ctxSubI#1#2{\ctxSub{#1}{#2}_I}

% Qualifs
\def\qKernel{\textsc{Kernel}\xspace}
\def\qUser{\textsc{User}\xspace}
\def\vTainted#1{\widehat{\diamondsuit}~#1}
\def\eTaint#1{\diamondsuit~#1}

\def\ptrK#1{#1~*}
\def\ptrU#1{#1~@}

% Eval de phrase
\def\ph#1#2#3{#1 ⊩ #2 → #3}

\def\phstar#1#2#3{#1 ⊩ #2 →^* #3}

\def\phprog#1#2{⊩ #1 →^* #2}

% typage
\def\ty#1#2#3{#1 ⊢ #2 : #3}
\def\tyi#1#2{#1 ⊢ #2}
\def\typh#1#2#3{#1 ⊢ #2 → #3}
\def\typhstar#1#2#3{#1 ⊢ #2 →^* #3}

\def\sVal{\textsc{Val}}
\def\sByte{\textsc{Byte}}

\def\powerset#1{\mathcal P \left({#1}\right)}

% Functions with placeholders
\def\phx#1{\ensuremath{#1{\cdot}}}
\def\phxx#1{\ensuremath{#1{\cdot}{\cdot}}}

% Progs bien formés
\def\WF#1{\textsc{wf}(#1)}
\def\NORET#1{\textsc{NoRet}(#1)}

\newcommand{\insertcode}[1]{\input{gen/#1.pyg.tex}}

% \chapter* + add to toc
%\newcommand{\chapterstar}[1]{\chapter*{#1} \addcontentsline{toc}{chapter}{#1}}

% Register names
\def\defreg#1{\expandafter\def\csname#1\endcsname{\textsc{#1}\xspace}}

\defreg{eax}
\defreg{ebx}
\defreg{ecx}
\defreg{edx}
\defreg{esi}
\defreg{edi}
\defreg{ebp}
\defreg{esp}
\defreg{eip}
\defreg{eflags}

\def\crtrois{\textsc{CR3}}

\newcommand{\ring}[1]{\emph{ring}~#1}

\newcommand{\asminstr}[1]{\texttt{#1}}

\newcommand{\wip}[0]{%
\begin{center}
    \includegraphics[width=5cm]{gen/travaux.pdf}
\end{center}
%\textbf{(wip)}
}

\def\listcons{\texttt{(::)}\xspace}

\tikzstyle{bigbrace} = [decorate,decoration={brace,amplitude=3mm}]

% Debug tikz bounding box
\def\drawbb{\draw[thick] (current bounding box.south west) rectangle (current bounding box.north east);}

% SubFloat + verbatim
\makeatletter
\newbox\sf@box
\newenvironment{SubFloat}[2][]%
{\def\sf@one{#1}%
\def\sf@two{#2}%
\setbox\sf@box\hbox
\bgroup}%
{ \egroup
\ifx\@empty\sf@two\@empty\relax
\def\sf@two{\@empty}
\fi
\ifx\@empty\sf@one\@empty\relax
\subfloat[\sf@two]{\box\sf@box}%
\else
\subfloat[\sf@one][\sf@two]{\box\sf@box}%
\fi}
\makeatother

% comes from the memoir manual {{{

\newcommand*{\addtotextheight}[1]{%
\twocolumn
\addtolength{\textheight}{#1}%
\onecolumn}

%% for main body, bottom of text at 1in, footer below
%% top of header at 1in, first text line double spaced
%% below base of header
\newlength{\linespace}
\setlength{\linespace}{\baselineskip} % current equivalent of \onelineskip
\setlength{\headheight}{\onelineskip}
\setlength{\headsep}{\linespace}
\addtolength{\headsep}{-\topskip}
\setlength{\uppermargin}{1in}
\addtolength{\uppermargin}{\headheight}
\addtolength{\uppermargin}{\headsep}

%% and for the bottom
\setlength{\lowermargin}{1in}
\setlength{\textheight}{\paperheight}

\addtolength{\textheight}{-\uppermargin}
\addtolength{\textheight}{-\lowermargin}
%% footnote settings
\setlength{\footskip}{\onelineskip}
\setlength{\footnotesep}{\onelineskip}


%% the fiddle lengths (..ta.. for title/approval page, others for prelims)
\newlength{\toptafiddle} \setlength{\toptafiddle}{2\linespace}
\newlength{\bottafiddle} \setlength{\bottafiddle}{0pt}
\newlength{\topfiddle}
\setlength{\topfiddle}{\toptafiddle}
\newlength{\botfiddle}
\setlength{\botfiddle}{\onelineskip}

%%% general macro for Abstract, etc., headings
\newcommand*{\pretoctitle}[1]{{\clearpage\centering \vspace*{-\topfiddle}#1\par}}

%%% Start the ACKNOWLEDGEMENTS
\newcommand{\asuacknowledgements}{\pretoctitle{REMERCIEMENTS}}

%%% Start the ABSTRACT
\newcommand{\asuabstract}{ \addtotextheight{-\botfiddle} \pretoctitle{ABSTRACT}}

\addtodef{\mainmatter}{\addtotextheight{\botfiddle}{}}

%% make it easy to center any dedication
\newcommand{\asudedication}[1]{%
{\clearpage\mbox{}\vfill\centering #1 \par\vfill\clearpage}}

% }}} /memoir

% {{{ subtitle command
\newcommand\subtitle[1]{%
  \newcommand\thesubtitle{#1}}
% }}}

% Language name
\def\langname{\textsc{Safespeak}\xspace}

% Center vertically + remove header & footer
\newenvironment{headingpage}{
\thispagestyle{simple}
\vspace*{\fill}
}{
\vspace*{\fill}
}

% Hoare triple
\def\hoare#1#2#3{\ensuremath{\{#1\}~#2~\{#3\}}}

\def\eqdef{\overset{\mathrm{def}}{=\joinrel=}\xspace}

% Lens
\def\mkLens#1#2{\langle #1 | #2 \rangle}
\def\setLens#1#2{\textsc{Lens}_{#1, #2}}
\def\lensGetX#1{\ensuremath{\mathrm{get}_{#1}}}
\def\lensPutX#1{\ensuremath{\mathrm{put}_{#1}}}
\def\lensGet#1#2{\lensGetX{#1}(#2)}
\def\lensPut#1#2#3{\lensPutX{#1}(#2, #3)}
\def\lensNodeBig#1#2{
\vcenter{\hbox{
\begin{tikzpicture}[scale=0.7]
\draw[fill=#1] (-0.5,-0.5) rectangle ++ (1,1);
\draw[fill=#2] (0,0) circle (0.3);
\end{tikzpicture}
}}
}
\def\lensInner#1{
\vcenter{\hbox{
\begin{tikzpicture}[scale=0.7]
\draw[fill=#1] (0, 0) circle (0.3);
\end{tikzpicture}
}}
}
\def\lensNodeBigIdx#1#2#3#4#5{
\vcenter{\hbox{
\begin{tikzpicture}[scale=0.7]
  \draw[fill=#1] (-1.5,-1.5) rectangle ++ (3,3);
  \node[fill=#2, draw, star, minimum size=6mm] at (-0.75, 0.75) {};
  \node[fill=#3, draw, star, minimum size=6mm] at ( 0.75, 0.75) {};
  \node[fill=#4, draw, circle, minimum size=6mm] at (-0.75,-0.75) {};
  \node[fill=#5, draw, regular polygon, regular polygon sides=5, minimum size=6mm] at ( 0.75,-0.75) {};
  \node at (-2,  0.75) {a};
  \node at ( 2,  0.75) {b};
  \node at (-2, -0.75) {c};
  \node at ( 2, -0.75) {d};
\end{tikzpicture}
}}
}
\def\lensInnerStar#1{
\vcenter{\hbox{
\begin{tikzpicture}[scale=0.7]
\node[draw, fill=#1, star] {};
\end{tikzpicture}
}}
}

% Erreurs

\def\serr#1{
  \switch{#1}
    {taint}{Ω_{taint}}
    {array}{Ω_{array}}
    {ptr}{Ω_{ptr}}
    {div}{Ω_{div}}
    {field}{Ω_{field}}
    {var}{Ω_{var}}
    {END}
    % important : étendre la description dans le chapitre 4 et la figure de
    % liste
}

%\def\uGet#1#2{#1 ⇐_U #2}
%\def\uPut#1#2{#1 ⇒_U #2}
\def\cfuF{\textrm{copy\_from\_user}\xspace}
\def\ctuF{\textrm{copy\_to\_user}\xspace}
\def\uGet#1#2{\cfuF(#1, #2)}
\def\uPut#1#2{\ctuF(#1, #2)}

% Litteral notation

\def\eStruct#1{\{#1\}}
\def\eArray#1{[#1]}
\def\tStruct#1{\eStruct{#1}}
\def\eFun#1#2{\mathrm{fun} (#1) \{ #2 \}}

\def\gramorx{\hspace{3mm}\gramor\xspace}
\def\stFun#1{\textsc{Fun}_n}
\def\tComp#1#2{\mathrm{Comp}(#1, #2)}

% Box on top on rules fig
\def\ruleheader#1{
  \begin{minipage}{\linewidth}
    \raggedleft
    \fbox{#1}
  \end{minipage}
}

\def\termFst#1{\ensuremath{\textrm{fst}~#1}}
\def\termSnd#1{\ensuremath{\textrm{snd}~#1}}

\def\ptrtype{\texttt{ptrtype}\xspace}

\def\addrInt{\textrm{addr}^\natural \xspace}
\def\interpPtr#1{\textrm{interp}_{(#1)}}

\tikzset{>=latex}

\def\partccl#1{
  \clearpage~\thispagestyle{cleared}\cleartorecto
  \chapter*{Conclusion de la partie~\ref{#1}}
  \addcontentsline{toc}{chapter}{Conclusion de la partie~\ref{#1}}
}

\def\hexa#1{0x#1}
\def\floatop#1{#1.}

\def\cLookup#1#2{\mathrm{Lookup} (#1, #2)}
\def\cExtend#1#2#3{\mathrm{Extend}(#1, #2 ↦ #3)}
\def\cCleanup#1{\mathrm{Cleanup}(#1)}

\newcommand{\motif}[1]{\cleaders\hbox{#1}\hfil}
\newcommand{\chain}[2]{\leavevmode\hbox to #2{\motif{#1}}}
\newcommand{\diamonds}{$\mkern-.6mu \diamond \mkern-.6mu$}

\def\jolibreak{
\fancybreak{\chain{\diamonds}{0.25\textwidth}}
}

\def\cDom#1{\mathrm{dom}(#1)}

\def\cVarname#1{\mathrm{Varname}(#1)}

\def\cVisible#1{\mathrm{Visible}(#1)}

\def\cRepr#1{\mathrm{Repr}(#1)}

\DeclareUnicodeCharacter{00D7}{\times}        % ×
\DeclareUnicodeCharacter{0393}{\Gamma}        % Γ
\DeclareUnicodeCharacter{03B5}{\epsilon}      % ε
\DeclareUnicodeCharacter{03BB}{\lambda}       % λ
\DeclareUnicodeCharacter{03C3}{\sigma}        % σ
\DeclareUnicodeCharacter{03C4}{\tau}          % τ
\DeclareUnicodeCharacter{2010}{-}             % ‐
\DeclareUnicodeCharacter{2026}{\ldots}        % …
\DeclareUnicodeCharacter{2190}{\leftarrow}    % ←
\DeclareUnicodeCharacter{2191}{\uparrow}      % ↑
\DeclareUnicodeCharacter{2192}{\rightarrow}   % →
\DeclareUnicodeCharacter{21A6}{\mapsto}       % ↦
\DeclareUnicodeCharacter{21D2}{\Rightarrow}   % ⇒
\DeclareUnicodeCharacter{2200}{\forall}       % ∀
\DeclareUnicodeCharacter{2205}{\emptyset}     % ∅
\DeclareUnicodeCharacter{2208}{\in}           % ∈
\DeclareUnicodeCharacter{2260}{\neq}          % ≠
\DeclareUnicodeCharacter{2264}{\le}           % ≤
\DeclareUnicodeCharacter{2265}{\ge}           % ≥
\DeclareUnicodeCharacter{2295}{\oplus}        % ⊕
\DeclareUnicodeCharacter{22A2}{\vdash}        % ⊢
\DeclareUnicodeCharacter{22D8}{\lll}          % ⋘
\DeclareUnicodeCharacter{22D9}{\ggg}          % ⋙


\usetikzlibrary{calc}
\usetikzlibrary{shapes}

\addtolength{\parskip}{\baselineskip}

\newcommand{\iIf}[3]{\textsc{if}~ #1 ~\textsc{then}~ #2 ~\textsc{else}~ #3}

\newcommand{\tInt}{\textsc{Int}}
\newcommand{\tFloat}{\textsc{Float}}
\newcommand{\tVoid}{\textsc{Void}}

\newcommand{\cNil}{\textsc{Nil}}

\newcommand{\ctonewspeak}{\textsc{c2newspeak}}
\newcommand{\cpp}{\textsc{cpp}}

\begin{document}

\section{Prétraitement}

Lors de l'écriture d'un compilateur, on a besoin d'un langage
intermédiaire qui fasse l'intermédiaire entre \emph{frontend} et
\emph{backend}. Avec ce langage, on doit pouvoir (facilement) réaliser
les types d'opérations suivantes :

\begin{itemize}
\item
  compiler ce langage vers un langage machine (interface avec le
  \emph{backend}).
\item
  exprimer des transformations intermédiaires sur cette représentation
  (analyses sémantiques, optimisations, etc).
\end{itemize}
À cause du premier point, un langage comme C est très séduisant, mais
malheureusement il a de nombreux défauts qui font qu'il ne remplit pas
le deuxième. Par exemple, certaines instructions sont ambiguës, et
l'ensemble des opération n'est pas orthogonal (on peut faire la même
chose de plusieurs manières).

Pour pallier ce problème, une première solution peut être d'ajouter des
contraintes au C intermédiaire manipulé pour en obtenir un sous-langage.
De nombreux sous-ensembles ont été définis pour aller dans ce sens.

\subsection{Critères}

Il y a une grande variété de langages intermédiaires différents : certains sont
plus adaptés comme support d'analyse statique, d'autres sont plutôt destinés à
générer du code dans un compilateur.

La question de l'interopérabilité est aussi importante pour l'ingéniérie :
certains langages intermédiaires se présentent sous la forme d'une bibliothèque
écrite dans un langage précis, alors que d'autres auront une représentation
textuelle.

Enfin, la maturité et la stabilité du langage intermédiaire doit être prise en
compte : sans documentation ou outils utilisables, un langage si parfait soit il
sera inutilisable.

\subsection{Candidats}

\paragraph{LLVM}\cite{llvm-pres} est un backend de compilateur développé par
Apple. Il se présente comme un concurrent direct de GCC et permet de compiler
\href{http://lwn.net/Articles/411654/}{le noyau Linux} et la majeure partie de
\href{http://lists.freebsd.org/pipermail/freebsd-current/2009-February/003743.html}{FreeBSD}
et \href{http://lwn.net/Articles/485301/}{Debian}.

LLVM est à la base une bibliothèque C++, mais une forme sérialisée sous forme
de fichier est également disponible.

\paragraph{Cminusminus}\cite{spjcmm} est un langage intermédiaire dont le but
est similaire à LLVM : pour quelqu'un qui écrit un compilateur, il est séduisant
de générer du C, mais pas pratique. À la place, autant générer du "C--".

\paragraph{Newspeak}\cite{npkSem} est un langage intermédiaire développé par
EADS dans le but d'être analysé par interprétation abstraite \todo{proofread}.

\paragraph{CMinor} est utilisé au sein du projet Compcert.

\paragraph{CIL}

\section{Syntaxe}

La grammaire suivante définit un langage impératif. On suppose qu'on peut
compiler un programme écrit en C vers ce langage.

Un programme est un triplet $P = (\vec{f}, \vec{x}, b)$ constitué d'un ensemble
de fonctions, d'un ensemble de variables et d'un bloc d'instructions. Ce bloc
sera exécuté au lancement du programme ; il peut par exemple contenir le code
d'initialisation des variables globales et l'appel à la fonction principale.

\begin{grammarf}
<programme> ::= ( fonctions , globales , <bloc> )
\end{grammarf}
\begin{minipage}{0.4\textwidth}
\begin{grammarf}
    <expr> ::= <lval>
          \alt unop <expr>
          \alt <expr> binop <expr>
          \alt cst
          \alt \& <lval>
          \alt \& fonction

    <lval> ::= var
          \alt <lval> . champ
          \alt <lval> [ <expr> ]
          \alt * <expr>
\end{grammarf}
\end{minipage}
\begin{minipage}{0.4\textwidth}
\begin{grammarf}
     <bloc> ::= <instr> ; <bloc>
           \alt $ε$

    <instr> ::= <lval> $←$ <expr>
           \alt <lval> $←$ funexp (args)
           \alt funexp (args)
           \alt $↑$ nom \{ <bloc> \}
           \alt if (<expr>) \{ <bloc> \} else \{ <bloc> \}
           \alt \{ <bloc> \} label:
           \alt goto label
           \alt forever \{ <bloc> \}
           \alt return <expr>
\end{grammarf}
\end{minipage}

\todo{expliquer pourquoi un while expr ne suffit pas}

Contrairement au langage C, le langage des expressions et des instructions est
séparé. Par conséquence, l'évaluation des expressions peut se faire sans effet de
bord.

\section{Compilation}

La compilation vers C est faite en trois étapes : prétraitement du code source,
compilation de C prétraité vers Newspeak, puis compilation de Newspeak vers ce
langage.

% {{{

\begin{tikzpicture}\shorthandoff{!}
\tikzstyle{file}=[draw, shape=rectangle, node distance=2.2cm, minimum
height=1cm, shade, top color=white,
    bottom color=blue!50!black!20, draw=blue!40!black!60, very thick];

\node [file] (c1) {\textcolor{black}{.c}};
\node [file, below of=c1] (c2) {\textcolor{black}{.c}};
\node [node distance=2.2cm, below of=c2] (c3) {};

\node [file, minimum height=0, node distance=5mm, above of=c3,draw] (c3b) {.adb};
\node [file, minimum height=0, node distance=5mm, below of=c3,draw] (c3s) {.ads};

\path (c3b.north west) ++(-3mm,3mm) [draw,dotted] rectangle ($(c3s.south east)+(3mm,-3mm)$);

\node [below of=c3, minimum height=3cm, node distance=2.2cm] (c4){};

\node [file, right of=c1] (cc1) {\textcolor{black}{.c}};
\node [file, right of=c2] (cc2) {\textcolor{black}{.c}};
\node [node distance=2.2cm, right of=c3] (cc3) {};

\node [below of=cc3, minimum height=3cm, node distance=2.2cm](cc4){};

\draw[->] (c1) -- node[above] {{\tiny \tt gcc -E}} (cc1);
\draw[->] (c2) -- (cc2);

\node [file, right of=cc1] (no1) {\textcolor{black}{.no}};
\node [file, right of=cc2] (no2) {\textcolor{black}{.no}};
\node [file, right of=cc3] (no3) {\textcolor{black}{.no}};

\node [below of=no3, minimum height=3cm, node distance=2.2cm](no4){};

\draw[->] (cc1) -- node[above] {{\tiny \tt c2newspeak -c}} (no1);
\draw[->] (cc2) --  (no2);

\draw[->] ($ (c3b.north east)!0.5!(c3s.south east) + (3mm,0) $) -- node[above] {\tiny \tt ada2newspeak -c} (no3);


\node [file, right of=no2] (npk) {\textcolor{black}{.npk}};

\node [right of=no3, minimum height=3cm, node distance=2.2cm](npk2){};
\node [right of=no4, minimum height=3cm, node distance=2.2cm](npk3){};

\node[draw, ellipse, right of=no2, minimum height=3cm]{};

\draw[->] (no1) -- node {} (npk);
\draw[->] (no2) -- node[above] {{\tiny \tt c2newspeak}} (npk);
\draw[->] (no3) -- node {} (npk);


\node [file, right of=npk](warn) {\textcolor{black}{.warn}};

\node [right of=npk3, node distance=3cm](analyser2){};

\draw[->] (npk) -- node[above] {{\tiny \tt penjili}} (warn);

\draw[->] (c4)   -- node[above, text depth=3pt] {\footnotesize Pr\'etraitement} (cc4);
\draw[->] (cc4)  -- node[above, text depth=3pt] {\footnotesize Compilation} (no4);
\draw[->] (no4)  -- node[above, text depth=3pt] {\footnotesize \'Edition de liens} (npk3);
\draw[->] (npk3) -- node[above, text depth=3pt] {\footnotesize Analyse} (analyser2);
\end{tikzpicture}

% }}}

\subsection{Prétraitement}

\ctonewspeak travaillant uniquement sur du code prétraité (dans directives de
préprocesseur), la première étape consiste donc à faire passer le code par \cpp.

\subsection{C vers Newspeak}

Cette passe est réalisée par l'utilitaire \ctonewspeak. L'essentiel de la
compilation consiste à mettre à plat les définition de types, et à simplifier le
flôt de contrôle. C en effet propose de nombreuses constructions ambigües ou
redondantes.

\subsection{Après Newspeak}

\todo{...}

\section{Sémantique}

On définit une sémantique opérationnelle à petits pas, sous forme d'une relation
de transition (notée $→$) entre états de l'interpréteur.

Ces états sont constitués des composantes suivantes :

\begin{itemize}
\item
  un point de contrôle $l$ dans le programme. Ils sont issus d'une
  première transformation en un graphe de flot de contrôle.
\item
  un état mémoire $σ$ qui associe à chaque variable une valeur.
\item
  une pile d'appels $s$ qui contient l'historique des appels de fonction
  imbriqués et leurs valeurs de retour.
\end{itemize}

\subsection{Graphe de flot de contrôle}

Dans la syntaxe ci-dessus, on peut classifier les instructions en deux familles:
celles qui définissent le flot de contrôle (if, dowith, goto, forever) et celles
qui définissent le flot de données. Une première transformation va transformer
chaque fonction en son graphe de flot de contrôle, défini comme suit :

\begin{itemize}
\item
  les nœuds sont des points de contrôle, qui représentent par exemple
  l'adresse mémoire de l'instruction qui vient d'être exécutée.
\item
  les arêtes sont soit des instructions "de données" (affectation,
  appel de fonction, déclaration), soit des conditions (ie une
  expression).
\end{itemize}

Intuitivement, on peut "passer" d'un état à un autre soit en passant par une
arête "condition" qui s'évalue à une valeur "vrai", soit en appliquant les
effets de bord d'une arête "instruction".

Dans la suite, on suppose qu'on a à notre disposition un ensemble de jugements :
$\langle l, instr, l' \rangle$ qui signifie qu'on peut passer du point $l$ au
point $l'$ en effectuant l'instruction $instr$.

\subsection{État mémoire}

La mémoire interne de l'interpréteur est une correspondance entre l'ensemble des
adresses (infini dénombrable) et l'ensemble des valeurs. Un état mémoire $σ$ est
une fonction partielle de $Addr$ vers $Val$.

Ces valeurs peuvent être de plusieurs formes :

\begin{align*}
v   ::=~ & n, f, \cNil & \text{constantes}\\
      |~ & \&a         & \text{pointeur vers l'adresse $a$}\\
      |~ & \&f         & \text{pointeur vers la fonction $f$}\\
      |~ & \top        & \text{valeur non initialisée}
\end{align*}

\todo{Pile d'appels}

\subsection{Left values}

La mémoire est organisée en adresses, mais pourtant dans le programme cette
notion n'est pas directement visible. Les accès sont réalisés à travers des
"left values".

\subsection{Jugements}

Les jugements ont les formes suivantes :

\begin{itemize}

\item $σ ⊢ lv ⇒ a$ :
  la left-value lv correspond à l'adresse mémoire a.

\item $σ ⊢ e ⇒ v$ :
  l'expression e s'évalue en v.

\item $(l, σ, s) → (l', σ', s')$ :
  permet de définir la fonction de transition principale

\end{itemize}

\subsection{Sémantique des left-values}

\begin{eqnarray*}
\irule{Eval-Lv-Var}{
  (v, a) ∈ σ
}{
  σ ⊢ v ⇒ a
}
&
\irule{Eval-Lv-Deref}{
  σ ⊢ e ⇒ \&a
}{
  σ ⊢ *e ⇒ a
}
\\
\irule{Eval-Lv-Field}{
  σ ⊢ lv ⇒ a
}{
  σ ⊢ lv.f ⇒ a + f
}
&
\irule{Eval-Lv-Array}{
  σ ⊢ lv ⇒ a \\
  σ ⊢ e ⇒ n \\
}{
  σ ⊢ lv[e] ⇒ a + n
}
\end{eqnarray*}

\subsection{Sémantique des expressions}

\begin{eqnarray*}
\irule{Eval-Cst}{
}{
  σ ⊢ c ⇒ c
}
&
\irule{Eval-Lv}{
  σ ⊢ lv ⇒ a \\
  (a, v) ∈ σ
}{
  σ ⊢ lv ⇒ v
}
\\
\irule{Eval-Unop}{
  σ ⊢ e ⇒ v
}{
  σ ⊢ op e ⇒ \widehat{op} v
}
&
\irule{Eval-Binop}{
  σ ⊢ e_1 ⇒ v_1 \\
  σ ⊢ e_2 ⇒ v_2
}{
  σ ⊢ e_1 op e_2 ⇒ v_1 \widehat{op} v_2
}
\\
\irule{Eval-AddrOf}{
  σ ⊢ lv ⇒ a
}{
  σ ⊢ \&lv ⇒ \&a
}
&
\irule{Eval-AddrOfFun}{
}{
  σ ⊢ \&f ⇒ \&f
}
\end{eqnarray*}

\subsection{Sémantique des instructions}

La règle la plus simple concerne l'affectation : on peut affecter une
expressions à une left value si elles ont le même type.

\[
\irule{Instr-Assign}{
  \langle l, lv ← e, l' \rangle \\
  σ ⊢ lv ⇒ a \\
  σ ⊢ e ⇒ v
}{
  (l, σ, s) → (l', σ [ a ↦ v ], s)
}
\]

Déclarer une variable, c'est rendre accessible dans un bloc une variable non
initialisée, qui n'est plus accessible par la suite : Si on suppose qu'on peut
traverser le bloc interne $b$ sous un $σ$ enrichi d'une nouvelle variable $x$,
on peut donc traverser l'instruction $↑x\{b\}$.

\begin{minipage}{0.6\textwidth}
\[
\irule{Instr-Decl}{
  \langle l, ↑x \{ b \}, l' \rangle \\
  \langle l_b, b, l_b' \rangle \\
  σ' = σ ⊕ \{ x → \top \} \\
  (l_b, σ', s) → (l_b', σ'', s')
}{
  (l, σ, s) → (l', σ'' \backslash x, s')
}
\]
\end{minipage}
\begin{minipage}{0.4\textwidth}
\begin{tikzpicture}
  [cfgnode/.style={draw,shape=circle},node distance=2cm]
  \node[cfgnode] (l) {$l$};
  \node[cfgnode,right of=l, node distance=4cm] (lp) {$l'$};
  \node[cfgnode,below of=l, xshift=1cm] (lb) {$l_b$};
  \node[cfgnode,below of=lp, xshift=-1cm] (lbp) {$l_b'$};
  \draw[->] (l)  to node[auto] {$↑x\{b\}$} (lp);
  \draw[->] (lb) to node[auto] {$b$} (lbp);
  \draw[->,dashed] (l) to node[auto] {$⊕ \{ x → \top \}$}(lb);
  \draw[->,dashed,swap] (lbp) to node[auto] {$\backslash x$} (lp);
\end{tikzpicture}
\end{minipage}

\todo{fcall}

%TODO pour :
%
%\[
%\irule{Instr-Fcall}{
%  \langle l, lv ← fe(\vec{e}), l' \rangle \\
%  σ ⊢ fe ⇒ f \\
%  σ ⊢ \vec{e} ⇒ \vec{v} \\
%  σ' = σ ⊕ \{args(f) = \vec{v}\} ⊕ \{ !ret → \top \} \\
%  (Entry(f), σ') → (Exit(f), σ'') \\
%  σ'' ⊢ !ret ⇒ v_{ret} \\
%  σ'' ⊢ lv ⇒ a
%}{
%  (l, σ, s) → (l', σ'' \backslash (args(f) \cup \{!ret\}) ⊕ \{ a → v_{ret}\}, ?)
%}
%\]

%TODO definir $σ ⊢ fe ⇒ f$

\subsection{Sémantique des conditions}

On utilise un encodage similaire à la déclaration. Tout d'abord, on évalue la
condition dans un contexte $σ$. Si elle s'évalue en un entier non nul, et qu'une
transition à travers le bloc $i_t$ est possible, alors on peut faire passer à
travers le "if".

\begin{minipage}{0.5\textwidth}
\begin{tikzpicture}
  [cfgnode/.style={draw,shape=circle}, node distance=1.5cm]
  \node[cfgnode] (l) {$l$};
  \node[cfgnode,right of=l, node distance=4.5cm] (lp) {$l'$};
  \node[cfgnode,below of=l, xshift=5mm] (lt) {$l_t$};
  \node[cfgnode,below of=lp, xshift=-5mm] (ltp) {$l_t'$};
  \node[cfgnode,above of=l, xshift=5mm] (lf) {$l_f$};
  \node[cfgnode,above of=lp, xshift=-5mm] (lfp) {$l_f'$};
  \draw[->] (l)  to node[auto] {$\iIf{e}{i_t}{i_f}$} (lp);
  \draw[->] (lt) to node[auto] {$i_t$} (ltp);
  \draw[->] (lf) to node[auto] {$i_f$} (lfp);
  \draw[->,dashed] (l) to node [auto,swap] {$≠0$} (lt);
  \draw[->,dashed] (ltp) to (lp);
  \draw[->,dashed] (l) to node [auto] {$0$} (lf);
  \draw[->,dashed] (lfp) to (lp);
\end{tikzpicture}
\end{minipage}
\begin{minipage}{0.5\textwidth}
\[
\irule{If-True}{
  \langle l, \iIf{e}{i_t}{i_f}, l' \rangle \\
  σ ⊢ e ⇒ n \\
  n ≠ 0 \\
  \langle l_i, i_t, l_i' \rangle \\
  (l_i, σ, s) → (l_i', σ', s')
}{
  (l, σ) → (l', σ', s')
}
\]

\[
\irule{If-False}{
  \langle l, \iIf{e}{i_t}{i_f}, l' \rangle \\
  σ ⊢ e ⇒ 0 \\
  \langle l_i, i_f, l_i' \rangle \\
  (l_i, σ, s) → (l_i', σ', s')
}{
  (l, σ, s) → (l', σ', s')
}
\]
\end{minipage}

\section{Typage}

Dans cette section, on définit la notion de programme bien typé. L'analyse par
typage permet de vérifier qu'à chaque expression on peut associer un type, et ce
de manière cohérente entre plusieurs utilisations d'une variable.

Un jugement de typage est de la forme $Γ ⊢ e : τ$ et se lit "sous $Γ$, $e$ est
typable en $τ$". Un environnement de typage $Γ$ contient le contexte nécessaire
à l'analyse, c'est à dire le type des fonctions et variables du programme.

Les instructions et blocs, au contraire, n'ont pas de type. On note $Γ ⊢ i$ pour
"sous $Γ$, $i$ est bien typé", c'est à dire que ces sous expressions sont
typables et que les types sont en accord avec le flot de données (par exemple,
pour que l'instruction $lv ← e$ soit bien typée sous $Γ$, il faut que les types
de $lv$ et de $e$ puissent avoir le même type sous $Γ$).

Les types des valeurs sont :

\begin{align*}
τ   ::=~ & \tInt, \tFloat, \tVoid & \text{constantes}\\
    |~   &  a                     & \text{variable}\\
    |~   & (τ_1, …, τ_n) → τ_r    & \text{fonction}\\
    |~   & [τ]                    & \text{tableau}\\
    |~   & τ*                     & \text{pointeur}\\
    |~   & \{ f_1:τ_1
            ,    …
            , f_n:τ_n \}          & \text{structure}
\end{align*}

\subsection{Schémas de type}

On va associer à chaque variable globale un type. Mais faire de même pourrait
être trop restrictif. En effet, une fonction comme memcpy peut être utilisée
pour copier des tableaux d'entiers, mais aussi de flottants. On va donc associer
un schéma de types à chaque fonction.

\[
  σ ::= ∀ \vec{a} . τ
\]

En associant un schéma de type $σ$ à une fonction $f$, on indique que la
fonction pourra être utilisée avec tout type $τ$ qui est une instanciation de
$σ$.

\subsection{Programme}

Au niveau global, un programme P est bien typé (noté $⊢ P$) s'il existe un
environnement $Γ^0$ permettant de typer ses composantes (les fonctions, les
globales et le bloc d'initialisation).

\[
\irule{Prog}{
  Γ^0 = (\vec{σ}, \vec{τ}) \\
  Γ^0 ⊢ \vec{f} : \vec{σ} \\
  Γ^0 ⊢ \vec{x} : \vec{t} \\
  Γ^0 ⊢ b
}{
  ⊢ (\vec{f}, \vec{x}, b)
}
\]

\subsection{Flot de contrôle}

Les règles suivantes permettent de définir les jugements $Γ ⊢ i$. En résumé, les
instructions sont bien typées si leurs sous-instructions sont bien typées. La
seule règle supplémentaire concerne la condition du $if$ qui doit être typée en
\tInt.

\begin{eqnarray*}
&
\irule{Pass}{
}{
  Γ ⊢ ε
}
&
\irule{Seq}{
  Γ ⊢ s \\
  Γ ⊢ b
}{
  Γ ⊢ s;b
}
\\
&
\irule{Forever}{
  Γ ⊢ b
}{
  Γ ⊢ forever \{ b \}
}
&
\irule{If}{
  Γ ⊢ e : \tInt \\
  Γ ⊢ i_t \\
  Γ ⊢ i_f
}{
  Γ ⊢ \iIf{e}{i_t}{i_f}
} \\
&
\irule{Goto}{
}{
  Γ ⊢ goto l
}
&
\irule{DoWith}{
  Γ ⊢ b
}{
  Γ ⊢ \{ b \} l:
}
\end{eqnarray*}

\subsection{Left values}

On étend la relation de typage aux left values : chaque left value, vue comme
une expression, peut être typée.

\begin{eqnarray*}
\irule{Lv-Var}{
  (v, τ) ∈ Γ
}{
  Γ ⊢ v : τ
}
&
\irule{Lv-Field}{
  Γ ⊢ lv : τ_s \\
  (f, τ_f) ∈ τ_s
}{
  Γ ⊢ lv.f : τ_f
}
\\
\irule{Lv-Deref}{
  Γ ⊢ e : τ*
}{
  Γ ⊢ *e : τ
}
&
\irule{Lv-Array}{
  Γ ⊢ lv : [τ] \\
  Γ ⊢ e : \tInt
}{
  Γ ⊢ lv[e] : τ
}
\end{eqnarray*}

\subsection{Expressions}
\begin{eqnarray*}
\irule{Const-Int}{
}{
  Γ ⊢ n : \tInt
}
\\
\irule{Const-Float}{
}{
  Γ ⊢ f : \tFloat
}
\\
\irule{Const-Nil}{
}{
  Γ ⊢ \cNil : τ*
}
\\
\irule{Op-Int}{
  op ∈ \{+,-,×,/,\&,|,\&\&,||,⋘,⋙\} \\
  Γ ⊢ e_1 : \tInt \\
  Γ ⊢ e_2 : \tInt
}{
  Γ ⊢ e_1 op e_2 : \tInt
}
\\
\irule{Op-Float}{
  op ∈ \{+.,-.,×.,/.\} \\
  Γ ⊢ e_1 : \tFloat \\
  Γ ⊢ e_2 : \tFloat
}{
  Γ ⊢ e_1 \mathrm{op}~ e_2 : \tFloat
}
\\
\irule{Op-Cmp}{
  op ∈ \{=,≠,≤,≥,<,>\} \\
  Γ ⊢ e_1 : τ \\
  Γ ⊢ e_2 : τ \\
  τ ∈ \{ \tInt, \tFloat \}
}{
  Γ ⊢ e_1  \mathrm{op} e_2 : \tInt
}
\\
\irule{Unop-Minus}{
  τ ∈ \{ \tInt, \tFloat \} \\
  Γ ⊢ e : τ
}{
  Γ ⊢ -e : τ
}
\\
\irule{Unop-Not}{
  op ∈ \{ \tilde, ! \}
  Γ ⊢ e : \tInt
}{
  Γ ⊢ op e : \tInt
}
\\
\irule{AddrOf}{
  Γ ⊢ lv : τ
}{
  Γ ⊢ \&lv : τ*
}
\\
\irule{AddrOfFun}{
  Γ ⊢ f : σ \\
  τ ≤ σ
}{
  Γ ⊢ \&lv : τ
}
\end{eqnarray*}

\subsection{Fonctions}

Pour typer une fonction, on commence par ajouter ses paramètres dans
l'environnement de typage, et on type la définition de la fonction. Le type
résultant est généralisé.

\[
\irule{Fun}{
  Γ' = Γ ⊕ \{ args(f) = \vec{τ} \} \\
  Γ' ⊢ body(f) \\
  Γ' ⊢ !ret : τ_r
}{
  Γ ⊢ f : Gen(\vec{τ} → τ_r, Γ)
}
\]

\subsection{Instructions}

\begin{eqnarray*}
\irule{Decl}{
  Γ ⊕ \{x : τ\} ⊢ b
}{
  Γ ⊢ ↑x \{b\}
}
\\
\irule{Assign}{
  Γ ⊢ lv : τ \\
  Γ ⊢ e  : τ
}{
  Γ ⊢ lv ← e
}
\\
\irule{Fcall}{
  Γ ⊢ lv : τ_{ret} \\
  Γ ⊢ fe : σ \\
  Γ ⊢ \vec{e} : \vec{τ} \\
  (\vec{τ} → τ_r) ≤ σ
}{
  Γ ⊢ lv ← fe(\vec{e})
}
\end{eqnarray*}

\bibliographystyle{alpha}
\bibliography{typ}
\end{document}
